\documentclass[11pt,a4paper]{memoir}

\usepackage[utf8]{inputenc}
\usepackage[T1]{fontenc}
\usepackage{geometry}
\geometry{margin=2.5cm}
\usepackage{hyperref}
\usepackage{listings}
\usepackage{xcolor}
\usepackage{graphicx}
\usepackage{amsmath,amssymb}
\usepackage{booktabs}

\hypersetup{colorlinks=true,linkcolor=blue,urlcolor=blue}
\lstset{basicstyle=\ttfamily\small,breaklines=true,columns=fullflexible}

\title{\textbf{Distributed Systems Project:\\Permissioned Blockchain Inventory System}}
\author{\textbf{Cancian Matteo} \and \textbf{Del Pio Luogo Marco}}
\date{Version 2.0 --- \today}

\begin{document}
\maketitle

\begin{abstract}
This project implements a permissioned inventory ledger composed of a Flutter
client, four Python validator services, and four Elixir consensus processes.
Clients submit Ed25519-signed inventory transactions to one validator. Python
services validate and store transactions in a SQLite-backed mempool, while
Elixir nodes disseminate proposals and execute a leader-based PoA/BFT
protocol. A block is committed after a quorum of three validators has signed
the same block hash. The client distinguishes committed inventory from local
pending overlays and exposes transaction status, blockchain history, and
block signatures.
\end{abstract}

\tableofcontents
\clearpage

\chapter{System Overview}

\section{Goals}
The system provides an auditable inventory ledger for four known validator
nodes. Inventory movements are represented by signed transactions and are
never edited after commitment. An incorrect movement is corrected by a new
compensating movement, not by deleting the original record.

\section{Implemented Architecture}
Each logical validator consists of:
\begin{itemize}
  \item a Python/FastAPI service for authentication, transaction validation,
        SQLite persistence, inventory state, and the HTTP API;
  \item an Elixir service for validator gossip, proposer selection, and
        consensus message handling.
\end{itemize}
The Flutter application is a client only: it does not participate in
consensus and does not decide whether a transaction is committed.

\begin{table}[htbp]
\centering
\caption{Local deployment ports}
\begin{tabular}{lll}
\toprule
\textbf{Component} & \textbf{Container port} & \textbf{Local host mapping}\\
\midrule
Python API & 8000/HTTP & 8001--8004\\
Elixir cluster gossip & 4000/TCP & internal Docker network\\
Elixir HTTP bridge & 6000/HTTP & internal Docker network\\
\bottomrule
\end{tabular}
\end{table}

The demo intentionally uses HTTP rather than HTTPS. HTTPS, certificate
management, and encrypted inter-site transport are deployment hardening tasks,
not guarantees of the current demo.

\chapter{Transaction Protocol}

\section{Transaction Contract}
The client sends the following logical fields to
\texttt{POST /api/v1/chain/add\_transaction}:
\begin{lstlisting}
{
  "item_id": "...",
  "quantity": 2,
  "location": "...",
  "client_ID": "...",
  "client_public_key": "...",
  "trans_id": "...",
  "timestamp": "...",
  "item_catalog": {...},
  "signature": "..."
}
\end{lstlisting}
\texttt{quantity} is a non-zero signed integer: positive values add stock and
negative values remove stock. There is no \texttt{tx\_type}. The client signs
the canonical JSON payload containing the transaction fields, including its
public key; the transport-only \texttt{flag} is not signed.

The first Python node verifies the authenticated user's identity and the
client signature. It forwards a flat transaction to its Elixir counterpart
with \texttt{flag=transaction}. A transaction received from another validator
is validated and stored locally but is not forwarded again. The
\texttt{trans\_id} primary key makes repeated delivery idempotent.

\section{Mempool and Proposal Trigger}
Every validator stores accepted transactions as \texttt{pending}. Python
requests a proposal when the configured threshold is reached or when the
configured idle interval elapses while pending transactions exist. The
threshold and interval are operational settings; the Elixir round timeout is
not an independent block-creation timer.

The proposer is deterministic:
\[
  \mathrm{proposer} = \mathrm{validators}[
  (\mathrm{height}+\mathrm{round}) \bmod N].
\]
With four validators, the commit quorum is three distinct valid validator
votes. Each block contains at most the configured maximum number of
transactions.

\section{Consensus}
The proposer asks Python to construct and sign a block containing height,
round, previous hash, transactions, and block hash. The proposer signature
covers the tuple \texttt{(height, round, proposer\_id, block\_hash)}.

The Elixir nodes then execute:
\begin{enumerate}
  \item \textbf{PRE-PREPARE:} the proposer broadcasts the block and its
        signature;
  \item \textbf{PREPARE:} validators verify the proposal and exchange
        validator votes; the proposer requires a prepare quorum before
        opening the commit phase;
  \item \textbf{COMMIT:} validators sign the block hash and broadcast commit
        votes; a node writes the block only after receiving three valid
        distinct commit votes;
  \item \textbf{WRITE:} each Python validator persists the block, applies its
        inventory changes, and marks included transactions as
        \texttt{finalized}.
\end{enumerate}

Prepare votes are protocol-internal and are not persisted in the public block
record. The committed block stores the proposer signature and the commit
certificate:
\begin{lstlisting}
"proposer_id": "node_b",
"proposer_signature": "...",
"validator_signatures": {
  "commit": {
    "node_a": "...",
    "node_b": "...",
    "node_c": "..."
  }
}
\end{lstlisting}

If consensus cannot progress, a round timeout may produce signed
\texttt{ROUND\_CHANGE} messages. Round change is subject to a quorum and does
not erase a committed block. The current implementation does not guarantee
automatic progress if the selected proposer remains unavailable in every
round; this is a known limitation documented as future work.

\chapter{Client Behaviour}

\section{Committed State and Pending Overlay}
The authoritative inventory is reconstructed from committed blocks. After a
client submits a transaction, that client may see a clearly local pending
overlay while the transaction is awaiting commitment. Other clients see only
committed state.

The Flutter client keeps pending transactions in an in-memory cache. On every
inventory or transaction refresh it reconciles the cache with the authoritative
\texttt{/user\_transactions} response. A transaction whose status is
\texttt{finalized} or \texttt{rejected} is removed before the pending overlay
is applied. This prevents a committed delta from being applied twice and
prevents a stale pending card from remaining visible.

\section{Client Screens}
\begin{itemize}
  \item \textbf{Inventario}: committed quantities plus the requester's pending
        overlay;
  \item \textbf{Transazioni}: the authenticated user's transactions and their
        pending/finalized status;
  \item \textbf{Blockchain}: committed blocks, proposer identity, changes, and
        cryptographic signatures;
  \item \textbf{Profilo}: account information and aggregate statistics based
        on individual user transactions.
\end{itemize}

\chapter{Security and Fault Model}

Validator identities and Ed25519 public keys are configured before startup.
Private keys remain local to their validator. Client public-key trust is
established by the first Python validator for this demo; all peers then use
the public key carried by the transaction. Canonical JSON and Base64-encoded
Ed25519 signatures are used for signed protocol payloads.

The four-validator configuration is designed for at most one Byzantine
validator, with $N=3f+1$ and a quorum of $2f+1=3$. This quorum argument
protects safety under the stated assumptions; availability also depends on
the selected proposer and network connectivity. The demo uses HTTP and a
shared application token, so it must not be exposed as a production service
without transport and secret-management hardening.

\chapter{Validation}

\section{Automated Checks}
The repository includes Flutter widget tests, Python service tests, and
Elixir ExUnit tests. Formatting, static analysis, and container health checks
are part of the validation workflow.

\section{Load and Byzantine Tests}
\texttt{Server\_Python/load\_test.py} submits signed transactions through an
Elixir router and monitors committed blocks. Since a block is capped at 200
transactions, a run of 1,998 transactions requires about ten blocks. The
simulation endpoints can activate invalid signatures, split-brain,
equivocation, or leader-offline scenarios on a selected node. These modes are
for demonstrations and testing only.

\chapter{Limitations and Future Work}
\begin{itemize}
  \item validate transaction timestamps and reject transactions older than a
        configured maximum age;
  \item add an optional nonce/replay cache;
  \item persist the client pending cache when appropriate;
  \item add forwarding retry metrics, status dashboards, and alerts;
  \item formally specify safe proposer takeover and locked-round recovery;
  \item deploy HTTPS, secret management, and authenticated inter-site
        transport.
\end{itemize}

\chapter{Conclusion}
The implemented system separates client UX, Python state authority, and
Elixir consensus responsibilities. Transactions become visible globally only
after a three-validator commit certificate is persisted. The client may show
a local optimistic overlay, but it reconciles that overlay with committed
server state so that confirmation cannot duplicate inventory changes.

\end{document}